% =============================================================================
%  Curriculum Vitae — Caterina Pedri
%  Build:  make          (or: tectonic -X compile Caterina_Pedri_CV.tex)
% =============================================================================
\documentclass{pedricv}

% ------------------------------------------------------------ personal data ---
% Empty fields are not printed.
\cvset{firstname}{Caterina}
\cvset{lastname}{Pedri}
\cvset{title}{Fisioterapista}
\cvset{subtitle}{Muscoloscheletrico e Sport}
\cvset{availability}{Disponibile a iniziare immediatamente}
\cvset{photo}{assets/photo.jpg}

\cvset{email}{caterina.pedri@gmail.com}
\cvset{phone}{}                         % es. +39 000 000 0000
\cvset{location}{}                      % es. Trento, Italia
\cvset{linkedin}{caterina-pedri}        % slug after linkedin.com/in/
\cvset{website}{caterinapedri.com}

% Printed at the foot of the sidebar on the last page; remove to omit.
\cvprivacy{Autorizzo il trattamento dei dati personali contenuti nel presente
  curriculum vitae ai sensi del Regolamento UE 2016/679 (GDPR) e del
  D.Lgs.~196/2003 e successive modifiche.}

\cvmetadata{fisioterapia, fisioterapista, muscoloscheletrico, sport,
  riabilitazione, return to sport, strength and conditioning, personal trainer}

\begin{document}
\begin{paracol}{2}

% =============================================================================
%  SIDEBAR
% =============================================================================
\cvphoto

\cvhascontacts{%
  \cvsection{Contatti}
  \cvcontacts
}{}

\cvsection{Competenze}
\begin{cvitems}
  \item Valutazione e gestione \hl{evidence-based} delle problematiche
        muscoloscheletriche, ortopediche e sportive
  \item Riabilitazione post-infortunio e post-chirurgica
  \item Esercizio terapeutico e terapia manuale, quando indicata
  \item Recupero di \hl{forza, mobilità, controllo motorio e capacità fisica}
  \item Gestione del carico e programmazione dell’esercizio
  \item Return to Activity e Return to Sport
  \item Strength \& Conditioning
  \item Analisi del movimento e della biomeccanica
  \item Educazione del paziente e promozione dell’autonomia nella gestione
        del problema
  \item Approccio patient-centered ed evidence-based
  \item Lavoro in équipe interdisciplinare e multidisciplinare
\end{cvitems}

\cvsection{Lingue}
\cvlanguage{Italiano}{Madrelingua}
\cvlanguage{Inglese}{Fluente}
\vspace{1mm}
{\small\color{muted}\RaggedRight
  Disponibile a lavorare sia in italiano sia in inglese.\par}

\cvsection{Formazione}
\cvdegree{Laurea in Fisioterapia}
         {Università degli Studi di Verona}
         {Rovereto, Italia\enspace·\enspace 2018 – 2021}
         {110/110 e lode}

\cvsection{Certificazioni}
\cvcert{ART® – Active Release Techniques, Certified Provider}{Aprile 2026}
\cvcert{ACL Rehabilitation: From Surgery to Return to Play}{2024}
\cvcert{Strength \& Conditioning}{2022}
{\small\color{muted}\RaggedRight
  Aggiornamento professionale continuo attraverso la consultazione della
  letteratura scientifica e della ricerca in ambito riabilitativo e
  dell’esercizio.\par}

\cvsection{Pubblicazioni}
{\raggedright\hyphenpenalty=10000 \exhyphenpenalty=10000
  {\small\color{muted}Feller~D., \textcolor{ink}{\textbf{Pedri~C.}},
   Gozzer~P., Innocenti~T., Trentin~F.\par}
  \vspace{1mm}
  \textit{The Reporting of Somatic Sensory Training Interventions in
  Individuals After a Stroke Is Suboptimal: A Systematic Review and
  Meta-research Study}\par
  \vspace{1mm}
  {\small\color{accent}\textbf{American Journal of Physical Medicine
   \& Rehabilitation}}\par
  {\small\color{muted}102(8):\,701–706\enspace·\enspace 1 agosto 2023\par}
  \vspace{0.6mm}
  {\small\href{https://doi.org/10.1097/PHM.0000000000002188}
    {\color{accent}\faExternalLink*\enspace 10.1097/PHM.0000000000002188}}\par}

% =============================================================================
%  MAIN COLUMN
% =============================================================================
\switchcolumn

\cvheader

\cvsection{Profilo}
\begin{cvprose}
Sono Caterina, fisioterapista in ambito muscoloscheletrico e sportivo, atleta
e da sempre appassionata di sport e movimento.

Negli ultimi cinque anni ho lavorato come fisioterapista tra Italia e Paesi
Bassi, seguendo persone dalla popolazione generale agli sportivi. L’esperienza
ad Amsterdam mi ha permesso di confrontarmi con un contesto internazionale e di
ampliare il mio modo di intendere la riabilitazione.

Parallelamente, ad Amsterdam ho iniziato a lavorare come Personal Trainer in
studi high-end, esperienza che ha rafforzato la mia visione del continuum tra
riabilitazione, allenamento e salute.

Oggi integro fisioterapia, personal training e health coaching, mantenendo il
\hlaccent{ragionamento clinico come punto di partenza imprescindibile}. Sono
abituata a lavorare in contesti interprofessionali e multidisciplinari e
considero il confronto e la collaborazione tra professionisti parte integrante
della qualità del percorso di cura.
\end{cvprose}

\cvsection{Esperienza professionale}

\cventry{Fisioterapista}{Set 2025 – Set 2026}
        {Art of Physio}{Amsterdam, Paesi Bassi}
\cvdesc{Fisioterapista muscoloscheletrica, ortopedica e sportiva in uno studio
  privato ad accesso diretto, con pazienti dalla popolazione generale agli
  atleti.}
\begin{cvitems}
  \item Valutazione e gestione di problematiche muscoloscheletriche,
        ortopediche e sportive
  \item Integrazione del lavoro fisioterapico con \hl{gestione del carico,
        programmazione dell’esercizio e progressione dell’attività}
  \item Definizione di percorsi individualizzati di riabilitazione e recupero
        funzionale
  \item Progressione verso il \hl{return to activity e return to sport}, in
        relazione agli obiettivi e alle richieste specifiche della persona
\end{cvitems}

\cventry{Personal Trainer}{Set 2025 – Set 2026}
        {Healthy Bestari}{Amsterdam, Paesi Bassi}
\cvdesc{Personal Trainer in stretta collaborazione con lo studio di
  fisioterapia Art of Physio, dove lavoravo parallelamente come
  fisioterapista.}
\begin{cvitems}
  \item Progettazione e conduzione di programmi di allenamento personalizzati,
        integrando principi di riabilitazione
  \item Lavoro su \hl{forza, mobilità, performance e gestione del dolore}
  \item Integrazione di strategie relative a esercizio, nutrizione, sonno e
        recupero per supportare salute e performance nel lungo periodo
\end{cvitems}

\cventry{Personal Trainer}{Set 2025 – Gen 2026}
        {Omnia Personal Training}{Amsterdam, Paesi Bassi}
\cvdesc{Studio di personal training high-end rivolto principalmente a
  imprenditori e professionisti con programmi di vita e lavoro ad alta
  intensità.}
\begin{cvitems}
  \item Coaching finalizzato al miglioramento della performance fisica e
        mentale
  \item Progettazione di programmi di allenamento efficienti e orientati ai
        risultati, adattati alle esigenze e alla disponibilità individuale
  \item Consulenza strategica su nutrizione, integrazione e recupero
  \item Supporto allo sviluppo di resilienza, capacità fisica e sostenibilità
        dello stile di vita attraverso un approccio globale alla salute
\end{cvitems}

\cventry{Personal Trainer\enspace{\color{faint}·}\enspace Tirocinio}
        {Set 2025 – Ott 2025}
        {Momentum Training Club}{Amsterdam, Paesi Bassi}
\cvdesc{Studio di personal training high-end orientato a un approccio
  \hl{data-driven alla performance}.}
\begin{cvitems}
  \item Monitoraggio di nutrizione, recupero, sonno e fattori legati allo stile
        di vita
  \item Integrazione di principi di \hl{Strength \& Conditioning e
        fisioterapia} nella programmazione individualizzata
  \item Utilizzo di valutazioni periodiche e dati di performance per adattare
        progressivamente gli interventi
\end{cvitems}

\cventry{Fisioterapista}{Apr 2022 – Lug 2025}
        {Casa di Cura Eremo}{Arco, Italia}
\cvdesc{Riabilitazione motoria e neuromotoria di pazienti ortopedici,
  neurologici e cardiopolmonari.}
\begin{cvitems}
  \item Esercizio terapeutico e utilizzo di tecnologie riabilitative avanzate
  \item Terapia in acqua e training respiratorio
  \item Programmazione di percorsi progressivi di recupero e ritorno
        all’attività
\end{cvitems}

\cventry{Fisioterapista}{Mar 2022 – Apr 2022}
        {Ospedale San Pancrazio}{Italia}
\begin{cvitems}
  \item Riabilitazione di pazienti ortopedici, neurologici e cardiologici
  \item Riabilitazione del pavimento pelvico mediante tecniche di biofeedback
  \item Programmi di recupero funzionale per pazienti con dipendenze
\end{cvitems}

\cventry{Tirocini clinici}{2018 – 2021}
        {Diverse strutture ospedaliere e ambulatoriali}
        {Trentino-Alto Adige, Italia}
\cvdesc{Quasi un anno complessivo di tirocinio clinico in riabilitazione
  \hl{ortopedica, neurologica e cardiopolmonare}.}

\cvsection{About me}
\begin{cvprose}
Sono Caterina, atleta e da sempre appassionata di sport e movimento. Dal tiro
con l’arco a livello agonistico al CrossFit, passando per sci, running,
powerlifting, windsurf, ciclismo e allenamento della forza, ho sperimentato in
prima persona come il movimento possa mettere alla prova, motivare e ispirare.

La mia mentalità da sportiva e la mia passione per il movimento influenzano
profondamente il mio modo di lavorare come fisioterapista.

Per me, la fisioterapia non è soltanto un trattamento a breve termine. È
un’opportunità per imparare a muoversi meglio, conoscere il proprio corpo e
sviluppare strategie per prendersene cura nel lungo periodo.

Vedo l’infortunio non soltanto come qualcosa da superare, ma anche come
un’opportunità per conoscersi meglio, comprendere il proprio corpo e imparare.
\hlaccent{La riabilitazione non finisce quando il dolore passa e recuperare non
significa semplicemente tornare a ciò che si era prima.} Significa poter uscire
dal percorso cambiati: con maggiore consapevolezza, autonomia e nuovi strumenti
per affrontare ciò che viene dopo.

Perché un corpo che attraversa un infortunio non è semplicemente un corpo da
riportare a com’era prima. È un corpo che ha attraversato un’esperienza, si è
adattato, è cambiato e continua a cambiare.

È qui che, per me, la riabilitazione diventa qualcosa di più del recupero di
una funzione: diventa un processo attraverso cui imparare a leggere i segnali
del proprio corpo, capire di cosa ha bisogno e costruire gli strumenti per
continuare a muoversi, allenarsi e vivere in salute anche oltre la fine della
riabilitazione.
\end{cvprose}

\end{paracol}
\end{document}
